\documentclass{beamer}

\mode<presentation>





\usetheme{Frankfurt}%
\usecolortheme{seagull}
\logo{\includegraphics[height=.25in]{clarksonGreen}}

\definecolor{garnet}{RGB}{136,0,0}
%\definecolor{clarksonGreen}{RGB}{0,71,28}
\definecolor{clarksonGreen}{RGB}{0,52,21}
\setbeamercolor{palette primary}{fg=clarksonGreen,bg=white}
\setbeamercolor{palette secondary}{fg=clarksonGreen,bg=white}
\setbeamercolor{palette tertiary}{fg=clarksonGreen,bg=white}
\setbeamercolor{palette quaternary}{bg=clarksonGreen,fg=white}
\setbeamercolor{block title}{fg=black,bg=black!15}
\setbeamercolor{block body}{fg=black,bg=black!10}
\setbeamercolor{titlelike}{bg=clarksonGreen,fg=white} % parent=palette quaternary}

\newcommand{\half}{\mbox{$\frac{1}{2}$}}
\newcommand{\deltat}{\mbox{$\triangle t$}}
\newcommand{\deltax}{\mbox{$\triangle x$}}
\newcommand{\deltay}{\mbox{$\triangle y$}}

\newcommand{\deriv}[2]{\frac{d}{d#2}#1}
\newcommand{\derivTwo}[2]{\frac{d^2}{d#2^2}#1}

\newcommand{\lp}{\left(}
\newcommand{\rp}{\right)}


\newcommand{\arrayTwo}[4]{
  \left[
  \begin{array}{rr}
    #1 & #2 \\
    #3 & #4
  \end{array}
  \right]
}

\newcommand{\vecTwo}[2]{
  \left[
  \begin{array}{r}
    #1 \\  #2
  \end{array}
  \right]
}


\newcommand{\stateTwo}[2]{
  \begin{array}{rr}
    \mbox{\fontsize{6}{6}\selectfont $#1$} \\  \mbox{\fontsize{6}{6}\selectfont $#2$}
  \end{array}
}


\newcommand{\arrayThree}[9]{
  \left[
    \begin{array}{rrr}
      #1 & #2 & #3 \\
      #4 & #5 & #6 \\
      #7 & #8 & #9
    \end{array}
  \right]
}

\newcommand{\startRowOpsTwo}{
  \left[
    \begin{array}{rr|rr}
}

\newcommand{\oneRowOpsTwo}[4] {
      #1 & #2 & #3 & #4 \\
}


\newcommand{\startRowOpsThree}{
  \left[
    \begin{array}{rrr|rrr}
}

\newcommand{\oneRowOpsThree}[6] {
      #1 & #2 & #3 & #4 & #5 & #6 \\
}

\newcommand{\stopRowOps}{
    \end{array}
  \right]
}


\newcommand{\vecThree}[3]{
  \left[
  \begin{array}{r}
    #1 \\  #2 \\ #3
  \end{array}
  \right]
}


\newcommand{\stateThree}[3]{
  \begin{array}{r}
    \mbox{\fontsize{6}{6}\selectfont $#1$} \\  
    \mbox{\fontsize{6}{6}\selectfont $#2$} \\ 
    \mbox{\fontsize{6}{6}\selectfont $#3$}
  \end{array}
}



\begin{document}

\title{Ordinary Differential Equations}
\subtitle{Math 232 - Week 6, Day 2}

\author{Kelly Black}
\institute{Clarkson University}
\date{5 Oct 2011}

\begin{frame}
  \titlepage
\end{frame}

\begin{frame}
  \frametitle{Outline}
  \tableofcontents[pausesection,hideallsubsections]
\end{frame}


\section{Matrix Inverse}


\begin{frame}
  \frametitle{Inverse of a Matrix}

  Suppose that we have the following system:
  \begin{eqnarray*}
    3x + y & = & 4 \\
    2x + y & = & -1
  \end{eqnarray*}

  Another way to express this:
  \begin{eqnarray*}
    \arrayTwo{3}{1}{2}{1} \vecTwo{x}{y} & = & \vecTwo{4}{-1}.
  \end{eqnarray*}

\end{frame}


\begin{frame}
  \frametitle{Interesting Coincidence}

  Note that
  \begin{eqnarray*}
    \arrayTwo{1}{-1}{-2}{3} \arrayTwo{3}{1}{2}{1}  & = & \arrayTwo{1}{0}{0}{1}.
  \end{eqnarray*}

  This means that we can go back to the previous equation to get
  \begin{eqnarray*}
    \arrayTwo{3}{1}{2}{1} \vecTwo{x}{y} & = & \vecTwo{4}{-1}, \\
    \arrayTwo{1}{-1}{-2}{3} \arrayTwo{3}{1}{2}{1} \vecTwo{x}{y} & = & 
          \arrayTwo{1}{-1}{-2}{3}\vecTwo{4}{-1}, \\
    \arrayTwo{1}{0}{0}{1} \vecTwo{x}{y} & = & \vecTwo{5}{-11}, \\
    \vecTwo{x}{y} & = & \vecTwo{5}{-11}.
  \end{eqnarray*}


\end{frame}

\section{Notation}

\begin{frame}
  \frametitle{Notation}

  \begin{eqnarray*}
    A & = & \arrayTwo{3}{1}{2}{1} \\
    B & = & \arrayTwo{1}{-1}{-2}{3} 
  \end{eqnarray*}

  Then
  \begin{eqnarray*}
    A\cdot B & = & I_3
  \end{eqnarray*}

\end{frame}


\begin{frame}
  \frametitle{Definition of the Inverse}

  If two matrices satisfy
  \begin{eqnarray*}
    A\cdot B & = & I
  \end{eqnarray*}
  then $B$ is the \textit{inverse} of $A$.

  It is denoted as $A^{-1}$.

  Note most matrices do not have an inverse! If $A^{-1}$ exists then
  we say that $A$ is \textit{invertible}.

\end{frame}

\section{Calculating the Inverse}

\begin{frame}
  \frametitle{Calculating the Inverse}

  How do we find $A^{-1}$?
  \begin{eqnarray*}
    \startRowOpsTwo
    \oneRowOpsTwo{3}{1}{1}{0}
    \oneRowOpsTwo{2}{1}{0}{1}
    \stopRowOps
  \end{eqnarray*}

  We put this in RREF. If we can make the left half look like the
  identity matrix then the right half is the inverse.

\end{frame}


\begin{frame}

  \begin{eqnarray*}
    \stateTwo{1/3 R_1}{~} 
    \startRowOpsTwo
    \oneRowOpsTwo{1}{1/3}{1/3}{0}
    \oneRowOpsTwo{2}{1}{0}{1}
    \stopRowOps \\
    \uncover<2->
    {
      \stateTwo{~}{-2R_1+R_2} 
      \startRowOpsTwo
      \oneRowOpsTwo{1}{1/3}{1/3}{0}
      \oneRowOpsTwo{0}{1/3}{-2/3}{1}
      \stopRowOps \\
    }
    \uncover<3->
    {
      \stateTwo{~}{3R_2} 
      \startRowOpsTwo
      \oneRowOpsTwo{1}{1/3}{1/3}{0}
      \oneRowOpsTwo{0}{1}{-2}{3}
      \stopRowOps 
    }
  \end{eqnarray*}

\end{frame}

\begin{frame}
  \begin{eqnarray*}
      \stateTwo{~}{3R_2} 
      \startRowOpsTwo
      \oneRowOpsTwo{1}{1/3}{1/3}{0}
      \oneRowOpsTwo{0}{1}{-2}{3}
      \stopRowOps \\
    \uncover<2->
    {
      \stateTwo{-1/3R_2+R_1}{~} 
      \startRowOpsTwo
      \oneRowOpsTwo{1}{0}{1}{-1}
      \oneRowOpsTwo{0}{1}{-2}{3}
      \stopRowOps 
    }
  \end{eqnarray*}

  \uncover<3->
  {
    The inverse of the matrix is 
    \begin{eqnarray*}
      A^{-1} & = & \arrayTwo{1}{-1}{-2}{3}
    \end{eqnarray*}
  }

\end{frame}


\begin{frame}
  \frametitle{Example}

  Is the matrix
  \begin{eqnarray*}
    \left[\begin{array}{rrr}
        2 & 4 & 4 \\
        1 & 1 & 2 \\
        0 & -1 & 2
      \end{array}\right]
  \end{eqnarray*}
  invertible?

  \uncover<2->
  {
    \begin{eqnarray*}
      \startRowOpsThree
      \oneRowOpsThree{2}{4}{4}{1}{0}{0}
      \oneRowOpsThree{1}{1}{2}{0}{1}{0}
      \oneRowOpsThree{0}{-1}{2}{0}{0}{1}
      \stopRowOps
    \end{eqnarray*}
  }

\end{frame}



\begin{frame}

    \begin{eqnarray*}
      \startRowOpsThree
      \oneRowOpsThree{2}{4}{4}{1}{0}{0}
      \oneRowOpsThree{1}{1}{2}{0}{1}{0}
      \oneRowOpsThree{0}{-1}{2}{0}{0}{1}
      \stopRowOps \\
      \uncover<2->
      {
        \stateThree{1/2 R_1}{~}{~}
        \startRowOpsThree
        \oneRowOpsThree{1}{2}{2}{1/2}{0}{0}
        \oneRowOpsThree{1}{1}{2}{0}{1}{0}
        \oneRowOpsThree{0}{-1}{2}{0}{0}{1}
        \stopRowOps \\
      }
      \uncover<3->
      {
        \stateThree{~}{-R_1+R_2}{~}
        \startRowOpsThree
        \oneRowOpsThree{1}{2}{2}{1/2}{0}{0}
        \oneRowOpsThree{0}{-1}{0}{-1/2}{1}{0}
        \oneRowOpsThree{0}{-1}{2}{0}{0}{1}
        \stopRowOps 
      }
    \end{eqnarray*}


\end{frame}



\begin{frame}
  \begin{eqnarray*}
    \stateThree{~}{~}{~}
    \startRowOpsThree
    \oneRowOpsThree{1}{2}{2}{1/2}{0}{0}
    \oneRowOpsThree{0}{-1}{0}{-1/2}{1}{0}
    \oneRowOpsThree{0}{-1}{2}{0}{0}{1}
    \stopRowOps \\
    \uncover<2->
    {
      \stateThree{~}{-R_2}{~}
      \startRowOpsThree
      \oneRowOpsThree{1}{2}{2}{1/2}{0}{0}
      \oneRowOpsThree{0}{1}{0}{1/2}{-1}{0}
      \oneRowOpsThree{0}{-1}{2}{0}{0}{1}
      \stopRowOps  \\
    }
    \uncover<3->
    {
      \stateThree{~}{R_2+R_3}{~}
      \startRowOpsThree
      \oneRowOpsThree{1}{2}{2}{1/2}{0}{0}
      \oneRowOpsThree{0}{1}{0}{1/2}{-1}{0}
      \oneRowOpsThree{0}{0}{2}{1/2}{-1}{1}
      \stopRowOps 
    }
  \end{eqnarray*}
\end{frame}




\begin{frame}
  \begin{eqnarray*}
    \stateThree{~}{~}{~}
      \startRowOpsThree
      \oneRowOpsThree{1}{2}{2}{1/2}{0}{0}
      \oneRowOpsThree{0}{1}{0}{1/2}{-1}{0}
      \oneRowOpsThree{0}{0}{2}{1/2}{-1}{1}
      \stopRowOps \\
    \uncover<2->
    {
      \stateThree{~}{~}{1/2 R_3}
      \startRowOpsThree
      \oneRowOpsThree{1}{2}{2}{1/2}{0}{0}
      \oneRowOpsThree{0}{1}{0}{1/2}{-1}{0}
      \oneRowOpsThree{0}{0}{1}{1/4}{-1/2}{1/2}
      \stopRowOps  \\
    }
    \uncover<3->
    {
      \stateThree{-2R_3+R_1}{~}{~}
      \startRowOpsThree
      \oneRowOpsThree{1}{2}{0}{0}{1}{-1}
      \oneRowOpsThree{0}{1}{0}{1/2}{-1}{0}
      \oneRowOpsThree{0}{0}{1}{1/4}{-1/2}{1/2}
      \stopRowOps 
    }
  \end{eqnarray*}
\end{frame}


\begin{frame}
  \begin{eqnarray*}
      \stateThree{~}{~}{~}
      \startRowOpsThree
      \oneRowOpsThree{1}{2}{0}{0}{1}{-1}
      \oneRowOpsThree{0}{1}{0}{1/2}{-1}{0}
      \oneRowOpsThree{0}{0}{1}{1/4}{-1/2}{1/2}
      \stopRowOps \\
    \uncover<2->
    {
      \stateThree{-2R_2+R_1}{~}{~}
      \startRowOpsThree
      \oneRowOpsThree{1}{0}{0}{-1}{3}{-1}
      \oneRowOpsThree{0}{1}{0}{1/2}{-1}{0}
      \oneRowOpsThree{0}{0}{1}{1/4}{-1/2}{1/2}
      \stopRowOps 
    }
  \end{eqnarray*}

  \uncover<3->
  {
    \begin{eqnarray*}
      A^{-1} & =  & \arrayThree{-1}{3}{-1}{1/2}{-1}{0}{1/4}{-1/2}{1/2}
    \end{eqnarray*}
  }

\end{frame}

\section{Algebraic Properties}

\begin{frame}
  \frametitle{Algebra Games}

  If you cannot make an identity matrix on the left hand side then the
  matrix is \textit{not invertible}.

  Note: Suppose that two matrices have inverses: \\
  \begin{tabular}{lll}
    $A$ & has inverse & $A^{-1}$ \\
    $B$ & has inverse & $B^{-1}$
  \end{tabular}

  Then:
  \begin{eqnarray*}
    A \cdot A^{-1} & = & I \\
    B \cdot B^{-1} & = & I
  \end{eqnarray*}
\end{frame}


\begin{frame}
  \frametitle{Inverse of the product}

  If that is the case then we can find the inverse of the product of
  the two matrices:
  \begin{eqnarray*}
    \lp B^{-1} A^{-1} \rp \lp A B \rp & = & B^{-1} \lp A^{-1} A \rp B \\
    & = & B^{-1} I B \\
    & = & B^{-1} B \\
    & = & I
  \end{eqnarray*}

  the inverse of $AB$ is $B^{-1}A^{-1}$.

\end{frame}


\begin{frame}
  \frametitle{Solving Linear Systems}

  Suppose we want to solve a linear system
  \begin{eqnarray*}
    A \vec{x} & = & \vec{b}.
  \end{eqnarray*}

  \textbf{If} $A^{-1}$ exists then 
  \begin{eqnarray*}
    A^{-1} A \vec{x} & = & A^{-1} \vec{b} \\
    \vec{x} & = & A^{-1} \vec{b}
  \end{eqnarray*}

  If the matrix is invertible then the homogeneous solution is the
  zero vector.

\end{frame}

\begin{frame}
  \frametitle{Why?}

  The homogeneous solution is a solution to the following system:
  \begin{eqnarray*}
    A \vec{x}_h & = & \vec{0}.
  \end{eqnarray*}

  \textbf{If} $A^{-1}$ exists then 
  \begin{eqnarray*}
    A^{-1} A \vec{x}_h & = & A^{-1} \vec{0}, \\
    \vec{x}_h & = & A^{-1} \vec{0}, \\
    & = & \vec{0}.
  \end{eqnarray*}

  If the inverse exists then the solution is unique.

\end{frame}




\end{document}

% LocalWords:  Clarkson pausesection hideallsubsections
